\section{Manuscript Structure}
\label{sec:manuscript_structure}

This manuscript is an account of the research conducted throughout this work, written for the reviewers and jury members who will evaluate it.
It is nevertheless intended to remain as accessible as possible to readers with a background in \gls{ai}.
To this end, the background chapters introduce and recall the concepts on which the subsequent chapters rely, allowing readers familiar with \gls{dl} or \gls{rl} to follow the manuscript without having to consult external material.

The manuscript is organized into a sequence of chapters, each addressing a specific topic related to the research questions introduced previously.
The chapters are divided into three categories, distinguished by color in Figure~\ref{fig:introduction_manuscript_map}:

\begin{figure}[t]
    \centering
    \resizebox{\textwidth}{!}{
        \begin{tikzpicture}[
            chapter/.style={rectangle, draw, rounded corners, minimum height=1.6cm,
                            minimum width=2.9cm, align=center, font=\small, text width=2.7cm},
            legend/.style={rectangle, draw, rounded corners, font=\small,
                           inner xsep=8pt, inner ysep=5pt},
            >={Stealth},
        ]

        % Chapter number (small, on top) followed by the chapter name in bold.
        \newcommand{\chapterbox}[2]{{\footnotesize Chapter~\ref{#1}}\\[3pt]\textbf{#2}}

        \node[chapter, draw=myblue,  fill=myblue!12]  (dl)   at (0, 0)    {\chapterbox{chap:deep_learning}{Deep\\Learning}};
        \node[chapter, draw=mygreen, fill=mygreen!12] (ls)   at (3.5, 0)  {\chapterbox{chap:latent_space}{Latent\\Space}};
        \node[chapter, draw=myblue,  fill=myblue!12]  (rl)   at (7, 0)    {\chapterbox{chap:reinforcement_learning}{Reinforcement\\Learning}};
        \node[chapter, draw=mygreen, fill=mygreen!12] (xrl)  at (10.5, 0) {\chapterbox{chap:explainability_reinforcement_learning}{Explainability\\in \gls{rl}}};
        \node[chapter, draw=myred,   fill=myred!12]   (car)  at (14, 0)   {\chapterbox{chap:clustering_agent_representation}{\gls{car}\\Method}};
        \node[chapter, draw=myred,   fill=myred!12]   (exp)  at (17.5, 0) {\chapterbox{chap:experiments}{Experiments}};

        % Reading order.
        \draw[->, thick] (dl)  -- (ls);
        \draw[->, thick] (ls)  -- (rl);
        \draw[->, thick] (rl)  -- (xrl);
        \draw[->, thick] (xrl) -- (car);
        \draw[->, thick] (car) -- (exp);

        % The method chapter draws on both survey chapters, routed above the
        % main chain at two different heights (absolute coordinates, so the
        % two elbowed paths stay clear of the intermediate boxes and of
        % each other) rather than through them.
        \draw[->, dashed, mygreen!60!black]
            (3.5, 0.8) -- (3.5, 1.8) -- (14.1, 1.8) -- (14.1, 0.8);
        \draw[->, dashed, mygreen!60!black]
            (10.5, 0.8) -- (10.5, 1.25) -- (13.8, 1.25) -- (13.8, 0.8);

        % Legend: each label sits inside its color swatch, sized to its text.
        \node[legend, draw=mygreen, fill=mygreen!12] (legend_sv) at (8.75, -1.7) {Survey};
        \node[legend, draw=myblue,  fill=myblue!12, left=0.8cm of legend_sv]  (legend_bg) {Technical background};
        \node[legend, draw=myred,   fill=myred!12,  right=0.8cm of legend_sv] (legend_me) {Method};

        \end{tikzpicture}
    }
    \caption{Structure of the manuscript.}
    \label{fig:introduction_manuscript_map}
\end{figure}


\begin{itemize}
\item[] \textcolor{myblue}{Technical background} chapters introduce fundamental concepts and established techniques required to understand the subsequent contributions.
\item[] \textcolor{mygreen}{Survey} chapters review and organize the existing literature around a specific research question.
The organization adopted in each of these surveys are original and therefore constitute contributions of this manuscript.
\item[] \textcolor{myred}{Method} chapters present \gls{car}, our \gls{me} method for \gls{rl}.
\end{itemize}

The chapters of the manuscript are organized as follows:

\begin{description}[style=nextline, leftmargin=1.5em, itemsep=0.5em]
\item[\textcolor{myblue}{Chapter~\ref{chap:deep_learning}}\enspace Deep Learning]
This chapter introduces the fundamental concepts of \gls{dl} and artificial \glspl{nn}.
It establishes the notation and key concepts used throughout the manuscript and provides the necessary background for understanding \glspl{nn}.

\item[\textcolor{mygreen}{Chapter~\ref{chap:latent_space}}\enspace Latent Space Analysis]
This chapter provides a survey of existing approaches for analyzing the \glspl{ls} learned by \glspl{nn}.
It reviews the different techniques used to characterize and interpret \glspl{lr}.
This chapter addresses Research Question~\ref{rq:latent_space}.

\item[\textcolor{myblue}{Chapter~\ref{chap:reinforcement_learning}}\enspace Reinforcement Learning]
This chapter introduces the fundamental concepts of \gls{rl}.
It presents the formal framework and introduces the main concepts, notation and algorithms required to study \gls{rl}.

\item[\textcolor{mygreen}{Chapter~\ref{chap:explainability_reinforcement_learning}}\enspace Explainability in Reinforcement Learning]
This chapter reviews existing approaches to explainability in \gls{rl}.
It provides a structured overview of the \gls{xrl} literature.
This chapter addresses Research Question~\ref{rq:taxonomy}.

\item[\textcolor{myred}{Chapter~\ref{chap:clustering_agent_representation}}\enspace The CAR Method]
This chapter presents our proposed approach to explainability in \gls{rl}.
The method builds on the analysis of \glspl{lr} and aims to identify the internal mechanisms underlying the behavior of \gls{rl} agents.
This chapter constitutes the methodological contribution of the manuscript and addresses Research Question~\ref{rq:mechanisms}.

\item[\textcolor{myred}{Chapter~\ref{chap:experiments}}\enspace Experimental Evaluation of the CAR Method]
This chapter presents the implementation and experimental evaluation of the proposed method.
It describes the experimental protocol and analyzes the results obtained with the \gls{car} method.
This chapter studies the relevance of the criteria used to evaluate the results produced by the \gls{car} method, addressing Research Question~\ref{rq:evaluation}.

\item[Conclusion]
The final chapter summarizes the main contributions of the manuscript, discusses the limitations of the proposed approaches, and presents possible directions for future research.

\end{description}

The organization of the manuscript also allows readers to follow different reading paths depending on their background and interests.
Readers already familiar with a given field may skip its corresponding background chapter and directly consult the chapters that build upon it.
For instance, readers familiar with \gls{dl} or \gls{rl} may skip Chapter~\ref{chap:deep_learning} or Chapter~\ref{chap:reinforcement_learning}, respectively.

The notation and acronyms used throughout the manuscript are collected in the glossary at the end of the document.

Having described the overall structure of the manuscript, we can now turn to its first technical chapter, which introduces the major developments that have shaped modern \gls{ai}: \gls{dl}.
